\chapter{Verificación con TLA\textsuperscript{+}}
\label{cap:tla}

Los caps.~\ref{cap:typegraph}--\ref{cap:transacciones} fijaron el modelo
algebraico de sotFS. Este capítulo explica cómo se modela ese sistema en
TLA\textsuperscript{+}, cómo se ejecuta TLC sobre los specs, cómo se
interpretan los resultados y, sobre todo, qué \emph{no} prueba este
método.

\section{La biblioteca de specs}

sotFS ships con seis specs TLA\textsuperscript{+} en \texttt{formal/}:

\begin{table}[h]
\centering
\renewcommand{\arraystretch}{1.2}
\begin{tabularx}{\linewidth}{l r X}
  \toprule
  \textbf{Spec} & \textbf{LOC} & \textbf{Propiedad} \\
  \midrule
  \citaTLA{sotfs\_graph.tla}                & 518 & Well-formedness del TypeGraph; los 7 invariantes (cap.~\ref{cap:typegraph}). \\
  \citaTLA{sotfs\_transactions.tla}         & 317 & GTXN atomicidad; \texttt{commit} $\Rightarrow$ \emph{visible} $\wedge$ invariantes; \texttt{rollback} $\Rightarrow$ restaurado. \\
  \citaTLA{sotfs\_capabilities.tla}         & 291 & Cap monotonicity, revoke transitivo, epoch invalida stale handles. \\
  \citaTLA{sotfs\_crash.tla}                & 287 & WAL-based crash recovery; cualquier prefijo del log replays a un estado válido. \\
  \citaTLA{sotfs\_crash\_refinement.tla}    & 395 & Refinement: el modelo concreto con crashes implementa el abstracto bajo stuttering. \\
  \citaTLA{sotfs\_curvature.tla}            & 301 & Curvatura estructural acotada bajo mutaciones adversariales. \\
  \bottomrule
\end{tabularx}
\caption{Los seis specs TLA\textsuperscript{+} de sotFS y la propiedad que
cada uno verifica.}
\label{tab:tla-specs}
\end{table}

Cada spec tiene tres configs (\texttt{*.cfg}, \texttt{*\_medium.cfg},
\texttt{*\_large.cfg}) para acotar el espacio de estados. La
configuración \texttt{small} es la que pasa CI; \texttt{medium} es la
del nightly; \texttt{large} se corre manualmente y puede excederse en
RAM dependiendo del host.

\section{Anatomía de un spec: \texttt{sotfs\_graph}}

Cualquier spec TLA\textsuperscript{+} tiene la misma estructura:
constantes, variables, init, next, fairness y propiedades.

\subsection*{Constantes y variables}

\begin{lstlisting}[language=tla, caption={Estado del spec
\texttt{sotfs\_graph}.}, label={lst:sotfs-graph-vars}]
CONSTANTS
    InodeIds,    DirIds,    BlockIds,
    Names,       MaxOps

VARIABLES
    inodes,           \* Set of active inode IDs
    dirs,             \* Set of active dir IDs
    blocks,           \* Set of active block IDs
    inodeType,        \* InodeId -> {"regular", "directory"}
    linkCount,        \* InodeId -> Nat
    containsEdges,    \* Set of <<dir, inode, name>> triples
    pointsToEdges,    \* Set of <<inode, block, offset>> triples
    dirInode,         \* DirId -> InodeId (the inode `.` points to)
    blockRefcount,    \* BlockId -> Nat
    ops               \* counter, bounds state space
\end{lstlisting}

Las \texttt{CONSTANTS} se instancian en el \texttt{.cfg}:

\begin{lstlisting}[language=tla, caption={\texttt{sotfs\_graph.cfg}: la
config \texttt{small}.}]
CONSTANTS
    InodeIds = {I1, I2, I3}
    DirIds   = {D1, D2}
    BlockIds = {B1, B2}
    Names    = {a, b}
    MaxOps   = 4
\end{lstlisting}

Esos cuatro inodes, dos dirs, dos bloques, dos nombres y traza máxima
de cuatro pasos producen un espacio de estados de $\sim$\,$10^5$
estados explorables en segundos. La config \texttt{medium} sube a
$\{I1\ldots I5\}$, $\{D1, D2, D3\}$, etc., y produce $\sim$\,$10^7$
estados.

\subsection*{Acciones (operaciones POSIX)}

Cada acción del spec es la versión TLA de una operación POSIX. Por
ejemplo, \texttt{LinkInode}:

\begin{lstlisting}[language=tla]
LinkInode(d, i, name) ==
    /\ d \in dirs
    /\ i \in inodes
    /\ ~\E e \in containsEdges :
          e[1] = d /\ e[3] = name   \* gluing condition
    /\ containsEdges' = containsEdges
                          \cup {<<d, i, name>>}
    /\ linkCount' = [linkCount EXCEPT
                       ![i] = @ + 1]
    /\ ops' = ops + 1
    /\ UNCHANGED <<inodes, dirs, blocks, ..., dirInode>>
\end{lstlisting}

La cláusula \texttt{\textasciitilde\textbackslash E e \textbackslash in
containsEdges : e[1] = d /\textbackslash e[3] = name} es la gluing
condition: ``no existe una arista \econtains{} desde \texttt{d} con
\texttt{name}''. Es la versión TLA del \eexist{} POSIX, pero \emph{no}
como un chequeo \emph{ad hoc}: emerge del estilo DPO del spec.

\subsection*{El predicado \texttt{Next}}

\begin{lstlisting}[language=tla]
Next ==
    \/ \E d \in dirs, i \in inodes, n \in Names : LinkInode(d, i, n)
    \/ \E d \in dirs, n \in Names : CreateFile(d, n)
    \/ \E d \in dirs, n \in Names : Mkdir(d, n)
    \/ \E d \in dirs, n \in Names : Unlink(d, n)
    \/ \E d \in dirs, n \in Names, n2 \in Names :
           Rename(d, n, n2)
    \/ ops >= MaxOps /\ UNCHANGED vars   \* stutter at bound
\end{lstlisting}

Cada disyunción es una acción posible en un paso. TLC explorará el
producto cartesiano de todas las maneras de instanciar los parámetros
existenciales.

\subsection*{Invariantes}

\begin{lstlisting}[language=tla]
INVARIANTS
    TypeInvariant
    LinkCountConsistent
    UniqueNamesPerDir
    DirHasSelfRef
    NoDanglingEdges
    BlockRefcountConsistent
    NoDirCycles
    ReachableHaveLinks
\end{lstlisting}

Estos son la versión TLA de los siete invariantes del
cap.~\ref{cap:typegraph}, más \texttt{ReachableHaveLinks} (todo inode
alcanzable desde la raíz tiene $\texttt{linkCount} \geq 1$). TLC evalúa
cada uno en \emph{cada estado alcanzable}; si alguno falla, se imprime
el contraejemplo (la traza que llegó al estado violatorio).

\section{Cómo correr TLC}

\begin{lstlisting}[language=shellcmd]
# Setup (una vez).
mkdir -p formal/lib
curl -L -o formal/lib/tla2tools.jar \
  https://github.com/tlaplus/tlaplus/releases/latest/\
download/tla2tools.jar

# Todos los specs, todos los tamaños.
just formal

# O directamente:
cd formal && bash run_tlc.sh

# Un solo tamaño.
bash run_tlc.sh small
bash run_tlc.sh medium

# Un solo spec, todos los tamaños.
bash run_tlc.sh graph
\end{lstlisting}

Salida: por cada combinación spec $\times$ config, un log en
\texttt{formal/tlc\_output/} y una tabla resumen en stdout.

\section{Cómo leer un contraejemplo de TLC}

Si una invariante falla, TLC imprime la \emph{traza mínima} que llevó al
estado violatorio:

\begin{lstlisting}[language=shellcmd]
Error: Invariant LinkCountConsistent is violated.
The behavior up to this point is:
  State 0:
    inodes = {I1}
    linkCount = (I1 :> 1)
    ...
  State 1: Action `Mkdir(D1, "a")`
    inodes = {I1, I2}
    linkCount = (I1 :> 1 @@ I2 :> 2)
    ...
  State 2: Action `Unlink(D1, "a")`
    inodes = {I1, I2}              <-- bug: I2 todavía vivo
    linkCount = (I1 :> 1 @@ I2 :> 1)  <-- pero linkCount=1
\end{lstlisting}

La traza tiene tres elementos clave: (i) el estado inicial; (ii) cada
acción aplicada con sus parámetros instanciados; (iii) el estado final
donde la invariante falla. El diagnóstico es directo: \texttt{Unlink}
no decrementó \texttt{linkCount} consistentemente con la remoción de
la arista \econtains.

\section{Crash refinement: el spec más sofisticado}

\citaTLA{sotfs\_crash\_refinement.tla} no es un spec con sus propios
estados nuevos: es un \emph{spec de refinement} que mapea el estado
del modelo concreto (con \texttt{Crash} y \texttt{Recover}) al estado
del modelo abstracto (sin crashes), y verifica el teorema:

\begin{lstlisting}[language=tla]
THEOREM Spec_Concrete => Spec_Abstract
\end{lstlisting}

El mapping (cap.~\ref{cap:transacciones}) lleva un estado concreto
\texttt{<<diskState, walState, memState, crashed>>} a un estado
abstracto que es o bien \texttt{memState} (si no crasheado) o bien
\texttt{ReplayWAL(walState, diskState)} (si crasheado). TLC verifica
exhaustivamente que toda traza concreta es traza abstracta.

\section{Status TLC v0.2.5}

Reporte oficial (\citaCommit{de} \texttt{formal/README.md}):
\textbf{14/14 PASS} en \texttt{small} y \texttt{medium}. La config
\texttt{large} puede agotar RAM dependiendo del host; el gating de CI
es \texttt{medium}.

\section{Lo que TLA\textsuperscript{+} no prueba}

\begin{itemize}
  \item No prueba el código Rust. La correspondencia spec $\to$ código
    es por inspección humana, no por extracción mecánica.
  \item No es prueba inductiva. TLC verifica modelos finitos por
    enumeración exhaustiva. Modelos más grandes podrían tener bugs no
    detectados.
  \item No verifica timing. Las propiedades de fairness y de
    \emph{eventually} usan el modelo lineal estándar; latencias reales
    no aparecen.
  \item No reemplaza a Coq. Las pruebas Coq del cap.~\ref{cap:coq}
    son inductivas y cubren modelos infinitos (universalmente
    cuantificados). TLC y Coq son complementarios.
\end{itemize}

\section*{Resumen del capítulo}

\begin{itemize}
  \item sotFS tiene seis specs TLA\textsuperscript{+} (graph,
    transactions, capabilities, crash, crash refinement, curvature),
    ejecutables con TLC bajo configs small / medium / large.
  \item Cada spec sigue la estructura
    \texttt{CONSTANTS}, \texttt{VARIABLES}, \texttt{Init}, \texttt{Next},
    \texttt{Spec}, \texttt{INVARIANTS}; las acciones son la versión TLA
    de las reglas DPO.
  \item Status v0.2.5: 14/14 PASS en small y medium.
  \item El crash refinement es un meta-spec que prueba que el modelo
    concreto refina al abstracto: \texttt{Spec\_Concrete \(\Rightarrow\)
    Spec\_Abstract}.
  \item Limitaciones: TLC explora modelos finitos; la correspondencia
    spec-código es por inspección.
\end{itemize}

\section*{Ejercicios}

\begin{ejercicio}[\dificultad{$\star$}]
\label{ej:tla-1}
Liste los seis specs TLA\textsuperscript{+} de sotFS con la propiedad
que cada uno verifica.
\end{ejercicio}

\begin{ejercicio}[\dificultad{$\star$}]
\label{ej:tla-2}
¿Cuáles son las constantes que se instancian en \texttt{sotfs\_graph.cfg}
en la config \texttt{small}?
\end{ejercicio}

\begin{ejercicio}[\dificultad{$\star\star$}]
\label{ej:tla-3}
Lea la acción \texttt{LinkInode} del listado~\ref{lst:sotfs-graph-vars}.
Identifique (a) la gluing condition, (b) la mutación, (c) las
variables \texttt{UNCHANGED}. ¿Qué invariante TLA podría romperse si
se omitiera el incremento de \texttt{linkCount}?
\end{ejercicio}

\begin{ejercicio}[\dificultad{$\star\star$}]
\label{ej:tla-4}
Diseñe una acción TLA \texttt{Rmdir(d)} para sotFS, con su gluing
condition explícita. Compare con la regla DPO del cap.~\ref{cap:dpo}.
\end{ejercicio}

\begin{ejercicio}[\dificultad{$\star\star$}]
\label{ej:tla-5}
Estime el espacio de estados de \citaTLA{sotfs\_graph.tla} con la
config \texttt{medium} (5 inodes, 3 dirs, 3 bloques, 3 nombres,
\texttt{MaxOps}=6). Pista: razonar sobre el número de estados
alcanzables de \texttt{containsEdges}.
\end{ejercicio}

\begin{ejercicio}[\dificultad{$\star\star$}]
\label{ej:tla-6}
Lea el último contraejemplo en \texttt{formal/tlc\_output/}
(si TLC nunca falló, simule uno: borre el incremento de \texttt{linkCount}
en \texttt{LinkInode} y vuelva a correr). Reproduzca el bug y explíquelo
en términos del cap.~\ref{cap:typegraph}.
\end{ejercicio}

\begin{ejercicio}[\dificultad{$\star\star$}]
\label{ej:tla-7}
\citaTLA{sotfs\_crash\_refinement.tla} prueba un teorema:
\texttt{Spec\_Concrete \(\Rightarrow\) Spec\_Abstract}. ¿Qué propiedad
formal del cap.~\ref{cap:transacciones} se sigue de ese teorema?
\end{ejercicio}

\begin{ejercicio}[\dificultad{$\star\star\star$}]
\label{ej:tla-8}
Extienda \citaTLA{sotfs\_capabilities.tla} con una acción
\texttt{TemporalRevoke(c, t\_end)} que suspende una cap hasta un
timestamp futuro. ¿Qué invariante nuevo hay que agregar? ¿Cambia
\texttt{MonotonicAttenuation}?
\end{ejercicio}
